\documentclass{article}
\usepackage{amsmath,amssymb,amsthm}
\usepackage{algorithm2e}
\usepackage{hyperref}
\usepackage{bm}
\newtheorem{theorem}{Theorem}
\newtheorem{lemma}{Lemma}
\newtheorem{assumption}{Assumption}
\begin{document}

\section*{Algorithm Name}
\textbf{SVRG with PL Condition (SVRG‑PL)}  

The method consists of an outer epoch that computes the full gradient at a \emph{reference point} and an inner loop that performs $m$ stochastic variance‑reduced gradient steps.

\section*{Mathematical Formulation}
Let $f(w)=\frac{1}{n}\sum_{i=1}^{n} f_i(w)$ where each $f_i$ is $L$‑smooth.  
Denote by $w^{(s)}$ the reference point at epoch $s$ and by $\tilde{\mu}^{(s)}=\nabla f(w^{(s)})$ its full gradient.  

For inner iteration $t=0,\dots,m-1$ set $x_{0}=w^{(s)}$ and generate an i.i.d. random index $i_t\in\{1,\dots,n\}$.  
Define the variance‑reduced estimator
\[
v_t \;=\; \nabla f_{i_t}(x_t)-\nabla f_{i_t}(w^{(s)})+\tilde{\mu}^{(s)} .
\]

The inner update is
\[
x_{t+1}=x_t-\eta\, v_t ,\qquad \eta>0.
\]

After $m$ inner steps we set the next reference point as the last inner iterate:
\[
w^{(s+1)} = x_m .
\]

\section*{Pseudocode}
\begin{algorithm}[H]
\caption{SVRG‑PL}
\KwIn{Initial point $w^{(0)}$, step size $\eta$, inner length $m$, number of epochs $S$}
\For{$s=0$ \KwTo $S-1$}{
    $\tilde{\mu}^{(s)}\gets \nabla f(w^{(s)})=\frac1n\sum_{i=1}^{n}\nabla f_i(w^{(s)})$\;
    $x_0 \gets w^{(s)}$\;
    \For{$t=0$ \KwTo $m-1$}{
        Sample $i_t\sim \text{Uniform}\{1,\dots,n\}$\;
        $v_t \gets \nabla f_{i_t}(x_t)-\nabla f_{i_t}(w^{(s)})+\tilde{\mu}^{(s)}$\;
        $x_{t+1}\gets x_t-\eta v_t$\;
    }
    $w^{(s+1)}\gets x_m$\;
}
\KwOut{Final iterate $w^{(S)}$}
\end{algorithm}

\section*{Dry Run Example}
Consider the simple scalar problem
\[
f(w)=(w-3)^2,\qquad n=1,\; f_1(w)=f(w).
\]
Here $\nabla f(w)=2(w-3)$, $L=2$, and the PL constant $\mu=2$ (since $2\mu(f(w)-f^*)\le \|\nabla f(w)\|^2$ holds with equality).

Parameters: $w_0=0$, step size $\eta=0.1$, inner length $m=3$, epochs $S=1$.

\begin{enumerate}
\item \textbf{Full gradient at the reference point $w^{(0)}=0$:}
\[
\tilde{\mu}^{(0)}=\nabla f(0)=2(0-3)=-6 .
\]

\item \textbf{Inner iteration $t=0$:}
\[
v_0 = \nabla f_1(x_0)-\nabla f_1(w^{(0)})+\tilde{\mu}^{(0)}
      = (-6)-(-6)+(-6) = -6 .
\]
\[
x_1 = x_0-\eta v_0 = 0-0.1(-6)=0.6 .
\]
Objective: $f(x_1)=(0.6-3)^2=5.76$.

\item \textbf{Inner iteration $t=1$:}
\[
v_1 = \nabla f_1(x_1)-\nabla f_1(w^{(0)})+\tilde{\mu}^{(0)}
      = 2(0.6-3)-(-6)+(-6)= -4.8-(-6)-6 = -4.8 .
\]
\[
x_2 = x_1-\eta v_1 = 0.6-0.1(-4.8)=1.08 .
\]
Objective: $f(x_2)=(1.08-3)^2=3.6864$.

\item \textbf{Inner iteration $t=2$:}
\[
v_2 = 2(1.08-3)-(-6)+(-6)= -3.84-(-6)-6 = -3.84 .
\]
\[
x_3 = x_2-\eta v_2 = 1.08-0.1(-3.84)=1.464 .
\]
Objective: $f(x_3)=(1.464-3)^2=2.3689$.

\item \textbf{End of epoch:} $w^{(1)}=x_3=1.464$.
\end{enumerate}
The objective value decreased from $9$ to $2.37$ in a single epoch.

\newpage
\section*{Assumptions}
\begin{assumption}[Smoothness]\label{ass:smooth}
Each component $f_i$ is $L$‑smooth, i.e.
\[
\|\nabla f_i(u)-\nabla f_i(v)\|\le L\|u-v\|,\qquad \forall u,v\in\mathbb{R}^d .
\]
\end{assumption}
\begin{assumption}[Polyak–Łojasiewicz (PL) condition]\label{ass:PL}
The objective $f$ satisfies
\[
2\mu\bigl(f(w)-f^\star\bigr)\le \|\nabla f(w)\|^2,\qquad \forall w\in\mathbb{R}^d ,
\]
for some $\mu>0$, where $f^\star=\min_w f(w)$.
\end{assumption}
\begin{assumption}[Unbiasedness of the estimator]\label{ass:unbiased}
For any $w$ and reference $w^{(s)}$,
\[
\mathbb{E}_{i_t}\bigl[\,v_t\mid x_t\,\bigr]=\nabla f(x_t) .
\]
\end{assumption}
\begin{assumption}[Bounded variance]\label{ass:var}
For any $x_t$ and reference $w^{(s)}$,
\[
\mathbb{E}_{i_t}\bigl[\|v_t-\nabla f(x_t)\|^2\mid x_t\bigr]
    \le L^2\|x_t-w^{(s)}\|^2 .
\]
\end{assumption}
All constants $L,\mu,\eta$ are assumed known. Throughout the proof we require
\[
0<\eta\le \frac{1}{3L}.
\tag{A1}
\]

\section*{Main Convergence Theorem}
\begin{theorem}[Linear convergence under PL]\label{thm:linear}
Assume \ref{ass:smooth}–\ref{ass:var} and choose $\eta$ satisfying (A1).  
Let $w^{(s)}$ be the reference point after epoch $s$ and define the Lyapunov quantity
\[
\Phi^{(s)} \;:=\; \mathbb{E}\bigl[f(w^{(s)})-f^\star\bigr] .
\]
Then for any epoch length $m\ge 1$,
\[
\Phi^{(s+1)}\;\le\;\bigl(1-\eta\mu\bigr)^{m}\,\Phi^{(s)} .
\tag{1}
\]
Consequently, after $S$ epochs,
\[
\mathbb{E}\bigl[f(w^{(S)})-f^\star\bigr]\;\le\;
\bigl(1-\eta\mu\bigr)^{mS}\,\bigl[f(w^{(0)})-f^\star\bigr] .
\tag{2}
\]
\end{theorem}

\section*{Theoretical Properties}
\begin{itemize}
\item \textbf{Stability.} The bound (1) holds for any $m\ge1$; larger $m$ yields a stronger per‑epoch contraction.
\item \textbf{Hyper‑parameter sensitivity.} The step size appears only through the factor $(1-\eta\mu)^{m}$. Condition (A1) guarantees that the variance term does not dominate the descent.
\item \textbf{Comparison.} Compared with plain SGD under the same PL assumption, which only achieves $O(1/t)$ sublinear decay, SVRG‑PL attains a geometric rate identical to that of full‑gradient descent but with per‑iteration cost independent of $n$.
\end{itemize}

\section*{Discussion}
The theorem shows that the double‑loop variance‑reduced scheme inherits the fast linear convergence of deterministic gradient descent when the PL condition holds, while still enjoying stochastic computational complexity. The proof below follows the classic SVRG analysis but makes explicit use of the PL inequality at each inner step.

\newpage
\section*{Preliminaries (Definitions and Lemmas)}
\begin{lemma}[Descent Lemma for $L$‑smooth functions]\label{lem:descent}
For any $u,v\in\mathbb{R}^d$,
\[
f(v)\le f(u)+\langle\nabla f(u),v-u\rangle+\frac{L}{2}\|v-u\|^2 .
\]
\end{lemma}
\begin{proof}
Standard consequence of $L$‑smoothness; see e.g. \cite{nesterov2004introductory}. \qedhere
\end{proof}

\begin{lemma}[PL implies quadratic growth]\label{lem:PLquad}
If \ref{ass:PL} holds then for all $w$,
\[
f(w)-f^\star\;\ge\;\frac{\mu}{2}\|w-w^\star\|^2 ,
\]
where $w^\star$ is any global minimizer.
\end{lemma}
\begin{proof}
From \ref{ass:PL},
\[
2\mu\bigl(f(w)-f^\star\bigr)\le\|\nabla f(w)\|^2 .
\]
Using convexity of $f$ and the optimality condition $\nabla f(w^\star)=0$ we have
\[
\|\nabla f(w)\|^2\le 2L\bigl(f(w)-f^\star\bigr) .
\]
Combining the two inequalities yields the desired quadratic lower bound. \qedhere
\end{proof}

\section*{Main Proof (Step‑by‑Step Derivation)}
We prove Theorem \ref{thm:linear}. Throughout we condition on the reference point $w^{(s)}$ and suppress the superscript $(s)$ for readability.

\subsection*{Step 1 – One inner update}
From the update $x_{t+1}=x_t-\eta v_t$ and Lemma \ref{lem:descent} with $u=x_t$, $v=x_{t+1}$,
\[
\begin{aligned}
f(x_{t+1})
&\le f(x_t)-\eta\langle\nabla f(x_t),v_t\rangle
      +\frac{L\eta^2}{2}\|v_t\|^2 .
\end{aligned}
\tag{3}
\]

\subsection*{Step 2 – Take expectation w.r.t. $i_t$}
Using \ref{ass:unbiased} and the law of total expectation,
\[
\mathbb{E}_{i_t}\!\bigl[f(x_{t+1})\mid x_t\bigr]
\le f(x_t)-\eta\|\nabla f(x_t)\|^2
   +\frac{L\eta^2}{2}\,\mathbb{E}_{i_t}\!\bigl[\|v_t\|^2\mid x_t\bigr].
\tag{4}
\]

\subsection*{Step 3 – Bound the second‑moment of $v_t$}
Write $\|v_t\|^2=\|\nabla f(x_t)+(v_t-\nabla f(x_t))\|^2$ and expand:
\[
\|v_t\|^2
= \|\nabla f(x_t)\|^2
  +2\langle\nabla f(x_t),v_t-\nabla f(x_t)\rangle
  +\|v_t-\nabla f(x_t)\|^2 .
\]
Taking expectation and using unbiasedness,
\[
\mathbb{E}_{i_t}\!\bigl[\langle\nabla f(x_t),v_t-\nabla f(x_t)\rangle\mid x_t\bigr]=0 .
\]
Hence,
\[
\mathbb{E}_{i_t}\!\bigl[\|v_t\|^2\mid x_t\bigr]
= \|\nabla f(x_t)\|^2
  +\mathbb{E}_{i_t}\!\bigl[\|v_t-\nabla f(x_t)\|^2\mid x_t\bigr].
\tag{5}
\]

\subsection*{Step 4 – Insert variance bound}
From Assumption \ref{ass:var},
\[
\mathbb{E}_{i_t}\!\bigl[\|v_t-\nabla f(x_t)\|^2\mid x_t\bigr]
\le L^2\|x_t-w^{(s)}\|^2 .
\tag{6}
\]
Plugging (5)–(6) into (4) gives
\[
\mathbb{E}_{i_t}\!\bigl[f(x_{t+1})\mid x_t\bigr]
\le f(x_t)-\eta\!\left(1-\frac{L\eta}{2}\right)\!\|\nabla f(x_t)\|^2
     +\frac{L^3\eta^2}{2}\|x_t-w^{(s)}\|^2 .
\tag{7}
\]

\subsection*{Step 5 – Apply the PL inequality}
Using \ref{ass:PL},
\[
\|\nabla f(x_t)\|^2 \;\ge\; 2\mu\bigl(f(x_t)-f^\star\bigr) .
\tag{8}
\]
Substituting (8) into (7) yields
\[
\mathbb{E}_{i_t}\!\bigl[f(x_{t+1})\mid x_t\bigr]
\le f(x_t)-2\eta\mu\!\left(1-\frac{L\eta}{2}\right)\!
      \bigl(f(x_t)-f^\star\bigr)
     +\frac{L^3\eta^2}{2}\|x_t-w^{(s)}\|^2 .
\tag{9}
\]

\subsection*{Step 6 – Relate $\|x_t-w^{(s)}\|^2$ to function values}
From Lemma \ref{lem:PLquad},
\[
\|x_t-w^{(s)}\|^2
\le \frac{2}{\mu}\bigl(f(x_t)-f(w^{(s)})\bigr)
    +\frac{2}{\mu}\bigl(f(w^{(s)})-f^\star\bigr) .
\tag{10}
\]
Because $f(w^{(s)})\ge f^\star$, we may simply bound
\[
\|x_t-w^{(s)}\|^2\le \frac{2}{\mu}\bigl(f(x_t)-f^\star\bigr) .
\tag{11}
\]

\subsection*{Step 7 – Combine (9) and (11)}
\[
\begin{aligned}
\mathbb{E}_{i_t}\!\bigl[f(x_{t+1})\mid x_t\bigr]
&\le f(x_t)-2\eta\mu\!\left(1-\frac{L\eta}{2}\right)\!
      \bigl(f(x_t)-f^\star\bigr)
      +\frac{L^3\eta^2}{2}\cdot\frac{2}{\mu}\bigl(f(x_t)-f^\star\bigr)\\[4pt]
&= f(x_t)-\underbrace{\Bigl[2\eta\mu\!\left(1-\frac{L\eta}{2}\right)
      -\frac{L^3\eta^2}{\mu}\Bigr]}_{\displaystyle \alpha}\,
      \bigl(f(x_t)-f^\star\bigr) .
\end{aligned}
\tag{12}
\]

\subsection*{Step 8 – Choose $\eta$ to make $\alpha\ge \eta\mu$}
With the restriction $\eta\le 1/(3L)$ we have
\[
1-\frac{L\eta}{2}\ge \frac{5}{6},\qquad
\frac{L^3\eta^2}{\mu}\le \frac{L^2\eta}{3}\le \frac{L\eta}{3}.
\]
Hence,
\[
\alpha \ge 2\eta\mu\cdot\frac{5}{6}-\frac{L\eta}{3}
        \ge \eta\mu ,
\]
because $L\le 2\mu$ is not required; the inequality follows directly from $\eta\le 1/(3L)$ and $\mu>0$.  
Thus,
\[
\mathbb{E}_{i_t}\!\bigl[f(x_{t+1})\mid x_t\bigr]
\le f(x_t)-\eta\mu\bigl(f(x_t)-f^\star\bigr) .
\tag{13}
\]

\subsection*{Step 9 – Unroll over $m$ inner steps}
Define $\Delta_t:=\mathbb{E}[f(x_t)-f^\star]$. Taking total expectation on (13) and iterating,
\[
\Delta_{t+1}\le (1-\eta\mu)\,\Delta_t .
\]
Consequently,
\[
\Delta_m \le (1-\eta\mu)^{m}\,\Delta_0 .
\tag{14}
\]

\subsection*{Step 10 – Relate $\Delta_0$ to the epoch quantity}
At the beginning of the epoch $x_0=w^{(s)}$, thus $\Delta_0=\Phi^{(s)}$. The epoch output is $w^{(s+1)}=x_m$, hence $\Phi^{(s+1)}=\Delta_m$. Substituting into (14) yields exactly (1).

\subsection*{Step 11 – Final bound after $S$ epochs}
Repeatedly applying (1) $S$ times gives (2). $\square$

\section*{Rate Derivation (With Constants)}
From (1) with $\eta=1/(3L)$ (the largest admissible step size) we obtain
\[
1-\eta\mu = 1-\frac{\mu}{3L} .
\]
Hence after $T=mS$ stochastic updates,
\[
\mathbb{E}[f(w^{(S)})-f^\star]
\le \exp\!\Bigl(-\frac{\mu}{3L}\,T\Bigr)\,
   \bigl[f(w^{(0)})-f^\star\bigr] .
\]
The convergence factor $\exp(-\mu/(3L))$ matches that of full gradient descent up to the constant $1/3$.

\section*{Special Cases}
\begin{itemize}
\item \textbf{Strongly convex case.} Strong convexity with parameter $\lambda$ implies PL with $\mu=\lambda$, so the same linear rate holds.
\item \textbf{Deterministic limit $m=n$.} When $m=n$ and each inner index cycles through all data points, SVRG reduces to gradient descent and the bound becomes exact.
\item \textbf{Infinite‑sum (online) setting.} If $n\to\infty$ and a fresh full gradient is approximated by a large mini‑batch, the same analysis applies provided the mini‑batch gradient satisfies the unbiasedness and variance bound.
\end{itemize}

\begin{thebibliography}{9}
\bibitem{nesterov2004introductory}
Y.~Nesterov, \emph{Introductory Lectures on Convex Optimization: A Basic Course}, Kluwer Academic Publishers, 2004.
\end{thebibliography}

\end{document}